% Part: first-order-logic
% Chapter: syntax-and-semantics
% Section: satisfaction

\documentclass[../../../include/open-logic-section]{subfiles}

\begin{document}

\olfileid{fol}{syn}{sat}

\olsection{ارضای \article{formula} \printtoken{S}{formula}
  در \article{structure} \printtoken{S}{structure}}

\begin{explain}
مفاهیم پایه‌ای که از یک سو عبارت‌هایی چون ترم‌ها و !!{formula}s و از سوی
دیگر !!{structure}s را به هم مربوط می‌کنند، \emph{!!{value}} یک ترم و
\emph{ارضای} !!a{formula} هستند. به‌طور غیرصوری، !!{value} یک ترم
!!a{element} از !!a{structure} است---اگر ترم فقط یک ثابت باشد، !!{value}
آن همان شیئی است که !!{structure} به ثابت نسبت داده است؛ و اگر با
!!{function}s ساخته شده باشد، !!{value} آن از !!{value}sی ثابت‌ها و
تابع‌هایی محاسبه می‌شود که به نمادهای تابعِ درون ترم نسبت داده شده‌اند.
!!^a{formula} در !!a{structure} \emph{ارضا می‌شود} هرگاه تعبیری که به
محمول‌ها داده شده است، آن !!{formula} را در دامنهٔ !!{structure} صادق
سازد. این مفهوم ارضا به‌صورت استقرایی مشخص می‌شود: تعیین !!{structure}
مستقیماً می‌گوید !!{formula}sی اتمی چه هنگام ارضا می‌شوند، و ارضای یک
!!{formula} مرکب را بسته به رابط اصلی یا سور آن و ارضاشدن یا نشدن
!!{subformula}sی بلافصلش تعریف می‌کنیم.

حالت سورها در اینجا قدری ظریف است، زیرا !!{subformula} بلافصل یک
!!{formula} سوربسته دارای !!{variable}ی آزاد است و !!{structure}s
!!{value}sی !!{variable}s را تعیین نمی‌کنند. برای رفع این دشواری،
\emph{گمارش‌های متغیر} را نیز معرفی می‌کنیم و ارضا را نه صرفاً نسبت به
!!a{structure}، بلکه نسبت به !!a{structure} همراه با !!a{variable}
گمارش تعریف می‌کنیم.
\end{explain}

\begin{defn}[گمارش متغیر]
یک \emph{گمارش متغیر}~$s$ برای !!a{structure}~$\Struct{M}$ تابعی است که
هر !!{variable} را به !!a{element} از~$\Domain M$ می‌نگارد؛ یعنی
$s\colon \Var \to \Domain M$.
\end{defn}

\begin{explain}
!!^a{structure} به هر !!{constant} یک !!a{value} نسبت می‌دهد و گمارش
متغیر به هر متغیر یک ارزش نسبت می‌دهد. اما می‌خواهیم ترم‌های ساخته‌شده از
آن‌ها نیز !!{element}sی !!{domain} را بنامند. برای این منظور، !!{value}
ترم‌ها را به‌صورت استقرایی تعریف می‌کنیم. ارزش !!{constant}s و متغیرها
همان است که !!{structure} یا گمارش متغیر تعیین می‌کند؛ ارزش ترم‌های
پیچیده‌تر نیز با استفاده از تابع‌هایی که !!{structure} به !!{function}s
نسبت می‌دهد، به‌صورت بازگشتی محاسبه می‌شود.
\end{explain}

\begin{defn}[!!^{value} ترم‌ها]
اگر $t$ ترمی از زبان~$\Lang L$، $\Struct M$ یک !!{structure}
برای~$\Lang L$، و $s$ یک گمارش !!a{variable} برای~$\Struct M$ باشد،
\emph{!!{value}}~$\Value{t}{M}[s]$ چنین تعریف می‌شود:
\begin{enumerate}
\item \indcase{t}{c}{$\Value{\indfrm}{M}[s] = \Assign{\indcomplex}{M}$.}
\item \indcase{t}{x}{$\Value{\indfrm}{M}[s] = s(\indcomplex)$.}
\item \indcase{t}{\Atom{f}{t_1, \ldots, t_n}}{
\[
\Value{\indfrm}{M}[s] = \Assign{f}{M}(\Value{t_1}{M}[s], \ldots,
\Value{t_n}{M}[s]).
\]}
\end{enumerate}
\end{defn}

\begin{defn}[واریانتِ $x$]
اگر $s$ گمارش !!a{variable} برای !!a{structure}~$\Struct M$ باشد، هر
گمارش !!{variable}~$s'$ برای~$\Struct M$ که حداکثر در مقداری که به~$x$
نسبت می‌دهد با~$s$ تفاوت داشته باشد، \emph{واریانتِ $x$} از~$s$ نامیده
می‌شود. اگر $s'$ واریانتِ $x$ از~$s$ باشد، می‌نویسیم
$\varAssign{s'}{s}{x}$.
\end{defn}

\begin{explain}
توجه کنید که واریانتِ $x$ از یک گمارش~$s$ \emph{لازم نیست} چیزی متفاوت
به~$x$ نسبت دهد. در واقع، هر گمارش واریانتِ $x$ از خودش به شمار می‌آید.
\end{explain}

\begin{defn}
  اگر $s$ گمارش !!a{variable} برای !!a{structure}~$\Struct M$ و
  $m \in \Domain{M}$ باشد، آنگاه گمارش~$\Subst{s}{m}{x}$ گمارش متغیری
  است که چنین تعریف می‌شود:
  \[\Subst{s}{m}{x}(y) = \begin{cases}
    m & \text{اگر } y \ident x\\
    s(y) & \text{در غیر این صورت}.
  \end{cases}\]
\end{defn}

به بیان دیگر، $\Subst{s}{m}{x}$ همان واریانتِ خاصِ $x$ از~$s$ است که
!!{element} دامنهٔ~$m$ را به~$x$ نسبت می‌دهد و به !!{variable}sی غیر
از~$x$ همان چیزهایی را نسبت می‌دهد که $s$ نسبت می‌دهد.

\begin{defn}[ارضا]
\ollabel{defn:satisfaction}
ارضای !!a{formula}~$!A$ در !!a{structure}~$\Struct M$ نسبت به گمارش
!!a{variable}~$s$، که با $\Sat{M}{!A}[s]$ نشان داده می‌شود، به‌صورت
بازگشتی چنین تعریف می‌شود. (برای بیان «نه $\Sat{M}{!A}[s]$» می‌نویسیم
$\Sat/{M}{!A}[s]$.)
\begin{enumerate}
\tagitem{prvFalse}{%
  \indcase{!A}{\lfalse}{$\Sat/{M}{\indfrm}[s]$.}}{}

\tagitem{prvTrue}{%
  \indcase{!A}{\ltrue}{$\Sat{M}{\indfrm}[s]$.}}{}

\item \indcase{!A}{\Atom{R}{t_1, \dots, t_n}}{$\Sat{M}{\indfrm}[s]$
  اگر و تنها اگر $\langle \Value{t_1}{M}[s], \dots, \Value{t_n}{M}[s] \rangle \in
  \Assign{R}{M}$.}

\item \indcase{!A}{\eq[t_1][t_2]}{$\Sat{M}{\indfrm}[s]$ اگر و تنها اگر
  $\Value{t_1}{M}[s] = \Value{t_2}{M}[s]$.}

\tagitem{prvNot}{%
  \indcase{!A}{\lnot !B}{$\Sat{M}{\indfrm}[s]$ اگر و تنها اگر
    $\Sat/{M}{!B}[s]$.}}{}

\tagitem{prvAnd}{%
  \indcase{!A}{(!B \land !C)}{$\Sat{M}{\indfrm}[s]$ اگر و تنها اگر
    $\Sat{M}{!B}[s]$ و $\Sat{M}{!C}[s]$.}}{}

\tagitem{prvOr}{%
  \indcase{!A}{(!B \lor !C)}{$\Sat{M}{\indfrm}[s]$ اگر و تنها اگر
    $\Sat{M}{!B}[s]$ یا $\Sat{M}{!C}[s]$ (یا هر دو).}}{}

\tagitem{prvIf}{%
  \indcase{!A}{(!B \lif !C)}{$\Sat{M}{\indfrm}[s]$ اگر و تنها اگر
    $\Sat/{M}{!B}[s]$ یا $\Sat{M}{!C}[s]$ (یا هر دو).}}{}

\tagitem{prvIff}{%
  \indcase{!A}{(!B \liff !C)}{$\Sat{M}{\indfrm}[s]$ اگر و تنها اگر یا
    هم $\Sat{M}{!B}[s]$ و هم $\Sat{M}{!C}[s]$، یا نه
    $\Sat{M}{!B}[s]$ و نه $\Sat{M}{!C}[s]$.}}{}

\tagitem{prvAll}{%
  \indcase{!A}{\lforall[x][!B]}{$\Sat{M}{\indfrm}[s]$ اگر و تنها اگر
    برای هر !!{element}~$m \in \Domain M$،
    $\Sat{M}{!B}[\Subst{s}{m}{x}]$.}}{}

\tagitem{prvEx}{%
  \indcase{!A}{\lexists[x][!B]}{$\Sat{M}{\indfrm}[s]$ اگر و تنها اگر
  برای دست‌کم یک !!{element}~$m \in \Domain M$،
  $\Sat{M}{!B}[\Subst{s}{m}{x}]$.}}{}
\end{enumerate}
\end{defn}

\begin{explain}
گمارش‌های متغیر در
\iftag{notprvEx,notprvAll}{بند آخر}{دو بند آخر} اهمیت دارند.\iftag{prvAll}{
نمی‌توانیم ارضای $\lforall[x][!B(x)]$ را با عبارت «برای هر $m \in
\Domain{M}$، $\Sat{M}{!B(m)}$» تعریف کنیم.}{}\iftag{prvEx}{ نمی‌توانیم
ارضای $\lexists[x][!B(x)]$ را با عبارت «برای دست‌کم یک $m \in
\Domain{M}$، $\Sat{M}{!B(m)}$» تعریف کنیم.}{} دلیل این است که اگر
$m \in \Domain M$ باشد، $m$ نمادی از زبان نیست؛ پس $!B(m)$ یک
!!a{formula} نیست (یعنی $\Subst{!B}{m}{x}$ تعریف نشده است). همچنین
نمی‌توانیم فرض کنیم !!{constant}s یا ترم‌هایی در اختیار داریم که هر
!!{element} از~$\Struct{M}$ را بنامند، زیرا در تعریف !!{structure}s هیچ
شرطی این امر را ایجاب نمی‌کند. در زبان استاندارد، مجموعهٔ !!{constant}s
!!{denumerable} است؛ پس اگر $\Domain{M}$ !!{enumerable} نباشد، حتی
!!{constant}s کافی برای نامیدن هر شیء وجود ندارد.

این مسئله را با معرفی گمارش‌های !!{variable} حل می‌کنیم؛ این گمارش‌ها
اجازه می‌دهند متغیرها را مستقیماً به !!{element}sی دامنه پیوند دهیم. سپس
به‌جای اینکه مثلاً بگوییم
\iftag{prvEx}{$\lexists[x][!B(x)]$}{$\lforall[x][!B(x)]$} در~$\Struct M$
ارضا می‌شود اگر و تنها اگر \iftag{prvEx}{برای دست‌کم یک $m \in
\Domain{M}$، $\Sat{M}{!B(m)}$}{برای هر $m \in \Domain{M}$،
$\Sat{M}{!B(m)}$}، می‌گوییم این فرمول در~$\Struct M$ \emph{نسبت
به}~$s$ ارضا می‌شود اگر و تنها اگر $!B(x)$ نسبت به~$\Subst{s}{m}{x}$
\iftag{prvEx}{برای دست‌کم یک}{برای هر} $m \in \Domain M$ ارضا شود.
\end{explain}

\begin{ex}
بگذارید $\Lang{L} = \{a, b, f, R\}$، که در آن $a$ و $b$ از
!!{constant}s، $f$ یک !!{function} دوجایگاهی و $R$ یک !!{predicate}
دوجایگاهی است. !!{structure}~$\Struct{M}$ زیر را در نظر بگیرید:
\begin{enumerate}
\item $\Domain M = \{1, 2, 3, 4\}$
\item $\Assign{a}{M} = 1$
\item $\Assign{b}{M} = 2$
\item $\Assign{f}{M}(x, y) = x+y$ اگر $x+y \le 3$، و در غیر این صورت $= 3$.
\item $\Assign{R}{M} = \{\tuple{1, 1}, \tuple{1, 2}, \tuple{2, 3}, \tuple{2, 4}\}$
\end{enumerate}
تابع $s(x) = 1$ که به هر !!{variable} مقدار $1 \in \Domain{M}$ را نسبت
می‌دهد، یک گمارش متغیر برای~$\Struct{M}$ است.

در این صورت
\begin{align*}
\Value{f(a,b)}{M}[s] & = \Assign{f}{M}(\Value{a}{M}[s], \Value{b}{M}[s]).
\intertext{چون $a$ و $b$ از !!{constant}s هستند، $\Value{a}{M}[s]
  = \Assign{a}{M} = 1$ و $\Value{b}{M}[s] = \Assign{b}{M} = 2$. پس}
\Value{f(a,b)}{M}[s] & = \Assign{f}{M}(1, 2) = 1+2 = 3.
\intertext{برای محاسبهٔ ارزش $f(f(a,b),a)$ باید عبارت زیر را در نظر بگیریم:}
\Value{f(f(a,b),a)}{M}[s] & = \Assign{f}{M}(\Value{f(a, b)}{M}[s],
\Value{a}{M}[s]) = \Assign{f}{M}(3, 1) = 3,
\intertext{زیرا $3+1 > 3$. چون $s(x) = 1$ و $\Value{x}{M}[s] =
  s(x)$، همچنین داریم}
\Value{f(f(a,b),x)}{M}[s] & = \Assign{f}{M}(\Value{f(a, b)}{M}[s],
\Value{x}{M}[s]) = \Assign{f}{M}(3, 1) = 3,
\end{align*}

یک !!{formula} اتمی~$R(t_1, t_2)$ هنگامی ارضا می‌شود که چندتاییِ
ارزش‌های آرگومان‌هایش، یعنی $\tuple{\Value{t_1}{M}[s],
  \Value{t_2}{M}[s]}$، !!a{element} از~$\Assign{R}{M}$ باشد. برای نمونه،
$\Sat{M}{R(b,f(a,b))}[s]$، زیرا $\tuple{\Value{b}{M},
  \Value{f(a,b)}{M}} = \tuple{2, 3} \in \Assign{R}{M}$؛ اما
$\Sat/{M}{R(x, f(a,b))}[s]$، زیرا $\tuple{1, 3} \notin \Assign{R}{M}$.

برای تعیین اینکه فرمول نااتمی~$!A$ ارضا می‌شود یا نه، بندی از تعریف
استقرایی را به کار می‌برید که به رابط اصلی آن مربوط است. برای نمونه، رابط
اصلی در $R(a, a) \lif (R(b, x) \lor R(x, b))$ همان~$\lif$ است، و
\begin{align*}
  & \Sat{M}{R(a, a) \lif (R(b, x) \lor R(x, b))}[s] \text{ اگر و تنها اگر }\\
  & \qquad
  \Sat/{M}{R(a,a)}[s] \text{ یا } \Sat{M}{R(b, x) \lor R(x, b)}[s]
  \intertext{چون $\Sat{M}{R(a,a)}[s]$ (زیرا $\tuple{1,1} \in
    \Assign{R}{M}$)، هنوز نمی‌توانیم پاسخ را تعیین کنیم و نخست باید
    دریابیم آیا $\Sat{M}{R(b, x) \lor R(x, b)}[s]$:}
  & \Sat{M}{R(b, x) \lor R(x, b)}[s] \text{ اگر و تنها اگر }\\
  & \qquad \Sat{M}{R(b,x)}[s] \text{ یا } \Sat{M}{R(x, b)}[s]
  \intertext{و این حالت برقرار است، زیرا $\Sat{M}{R(x, b)}[s]$
    (چون $\tuple{1,2} \in \Assign{R}{M}$).}
\end{align*}

به یاد آورید که واریانتِ $x$ از~$s$ گمارش متغیری است که حداکثر در مقداری
که به~$x$ نسبت می‌دهد با $s$ تفاوت دارد. برای هر !!{element}
از~$\Domain{M}$، یک واریانتِ $x$ از~$s$ وجود دارد:
\begin{align*}
  s_1 & = \Subst{s}{1}{x}, &
  s_2 & = \Subst{s}{2}{x},\\
  s_3 & = \Subst{s}{3}{x}, &
  s_4 & = \Subst{s}{4}{x}.
\end{align*}
پس برای نمونه، $s_2(x) = 2$ و برای همهٔ متغیرهای~$y$ غیر از~$x$،
$s_2(y) = s(y) = 1$. این‌ها همهٔ واریانت‌های $x$ از~$s$ برای
ساختار~$\Struct{M}$ هستند، زیرا $\Domain{M} = \{1, 2, 3, 4\}$. به‌ویژه
توجه کنید که $s_1 = s$ ($s$ همیشه واریانتِ $x$ از خودش است).

\iftag{prvEx}{برای تعیین اینکه !!{formula} با سور وجودی
  $\lexists[x][!A(x)]$ ارضا می‌شود یا نه، باید تعیین کنیم آیا
  $\Sat{M}{!A(x)}[\Subst{s}{m}{x}]$ برای دست‌کم یک $m \in \Domain M$
  برقرار است. پس
  \[
  \Sat{M}{\lexists[x][(R(b,x) \lor R(x,b))]}[s],
  \]
  زیرا $\Sat{M}{R(b,x) \lor R(x, b)}[\Subst{s}{1}{x}]$
  ($\Subst{s}{3}{x}$ نیز کارساز بود). اما
  \[
  \Sat/{M}{\lexists[x][(R(b,x) \land R(x,b))]}[s]
  \]
  زیرا هر $m \in \Domain{M}$ را که برگزینیم،
  $\Sat/{M}{R(b,x) \land R(x,b)}[\Subst{s}{m}{x}]$.}{}

\iftag{prvAll}{برای تعیین اینکه !!{formula} با سور کلی
  $\lforall[x][!A(x)]$ ارضا می‌شود یا نه، باید تعیین کنیم آیا
  $\Sat{M}{!A(x)}[\Subst{s}{m}{x}]$ برای همهٔ $m \in \Domain M$ برقرار
  است. پس
  \[
  \Sat{M}{\lforall[x][(R(x,a) \lif R(a,x))]}[s],
  \]
  زیرا $\Sat{M}{R(x,a) \lif R(a,x)}[\Subst{s}{m}{x}]$ برای همهٔ
  $m \in \Domain M$. برای $m = 1$، داریم
  $\Sat{M}{R(a,x)}[\Subst{s}{1}{x}]$، پس تالی صادق است؛ برای $m = 2$،
  $3$ و~$4$، داریم $\Sat/{M}{R(x,a)}[\Subst{s}{m}{x}]$، پس مقدم کاذب
  است. اما
  \[
  \Sat/{M}{\lforall[x][(R(a,x) \lif R(x,a))]}[s]
  \]
  زیرا $\Sat/{M}{R(a,x) \lif R(x,a)}[\Subst{s}{2}{x}]$ (چون
  $\Sat{M}{R(a, x)}[\Subst{s}{2}{x}]$ و $\Sat/{M}{R(x,
  a)}[\Subst{s}{2}{x}]$).}{}

\iftag{defEx}{برای تعیین اینکه !!{formula} با سور وجودی
  $\lexists[x][!A(x)]$ ارضا می‌شود یا نه، باید تعیین کنیم آیا
  $\Sat{M}{\lnot \lforall[x][\lnot !A(x)]}[s]$. برای نمونه، داریم
  \[
  \Sat{M}{\lexists[x][(R(b,x) \lor R(x,b))]}[s].
  \]
  نخست، $\Sat{M}{R(b,x) \lor R(x, b)}[\Subst{s}{1}{x}]$
  ($\Subst{s}{3}{x}$ نیز کارساز بود). پس
  $\Sat/{M}{\lnot(R(b,x) \lor R(x,b))}[\Subst{s}{1}{x}]$، بنابراین
  $\Sat/{M}{\lforall[x][\lnot((R(b,x) \lor R(x,b)))]}[s]$ و در نتیجه
  $\Sat{M}{\lnot\lforall[x][\lnot((R(b,x) \lor R(x,b)))]}[s]$. از سوی دیگر،
  \[
  \Sat/{M}{\lexists[x][(R(b,x) \land R(x,b))]}[s].
  \]
  دلیل این است که $\Sat{M}{\lforall[x][\lnot(R(b,x) \land
  R(x,b))]}[s]$، زیرا هیچ $m \in \Domain M$ای نیست که
  $\Sat{M}{R(b,x) \land R(x,b)}[\Subst{s}{m}{x}]$. چنان‌که احتمالاً از
  این مثال‌ها حدس می‌زنید، $\Sat{M}{\lexists[x][!A(x)]}[s]$ اگر و تنها
  اگر $\Sat{M}{!A(x)}[\Subst{s}{m}{x}]$ برای دست‌کم یک
  $m \in \Domain M$.}{}

\iftag{defAll}{برای تعیین اینکه !!{formula} با سور کلی
  $\lforall[x][!A(x)]$ ارضا می‌شود یا نه، باید تعیین کنیم آیا
  $\Sat{M}{\lnot\lexists[x][\lnot !A(x)]}[s]$. برای نمونه،
  \[
  \Sat{M}{\lforall[x][(R(x,a) \lif R(a,x))]}[s],
  \]
  نخست، $\Sat{M}{R(x,a) \lif R(a,x)}[\Subst{s}{m}{x}]$ برای همهٔ
  $m \in \Domain M$ ($\Sat{M}{R(a,x)}[\Subst{s}{1}{x}]$ و
  $\Sat/{M}{R(a,x)}[\Subst{s}{m}{x}]$ برای $m = 2$، $3$ یا~$4$).
  بنابراین، هیچ $m \in \Domain M$ای نیست که
  $\Sat{M}{\lnot(R(x,a) \lif R(a,x))}[\Subst{s}{m}{x}]$ و ازاین‌رو
  $\Sat/{M}{\lexists[x][\lnot(R(x,a) \lif R(a,x))]}[s]$. پس
  $\Sat{M}{\lnot\lexists[x][\lnot(R(x,a) \lif R(a,x))]}[s]$. از سوی
  دیگر،
  \[
  \Sat/{M}{\lforall[x][(R(a,x) \lif R(x,a))]}[s],
  \]
  زیرا $\Sat/{M}{R(a,x) \lif R(x,a)}[\Subst{s}{2}{x}]$؛ پس
  $\Sat{M}{\lexists[x][\lnot(R(a,x) \lif R(x,a))]}[s]$. چنان‌که
  احتمالاً از این مثال‌ها حدس می‌زنید،
  $\Sat{M}{\lforall[x][!A(x)]}[s]$ اگر و تنها اگر
  $\Sat{M}{!A(x)}[\Subst{s}{m}{x}]$ برای هر $m \in \Domain M$.}{}

برای حالتی پیچیده‌تر، فرمول زیر را در نظر بگیرید:
\[
\lforall[x][(R(a,x) \lif \lexists[y][R(x,y)])].
\]
چون $\Sat/{M}{R(a,x)}[\Subst{s}{3}{x}]$ و
$\Sat/{M}{R(a,x)}[\Subst{s}{4}{x}]$، حالت‌های جالبی که باید در آن‌ها
نگران تالی شرطی باشیم فقط $m = 1$ و $m = 2$ هستند. آیا
$\Sat{M}{\lexists[y][R(x,y)]}[\Subst{s}{1}{x}]$ برقرار است؟ این امر
هنگامی برقرار است که دست‌کم یک $n \in \Domain M$ وجود داشته باشد که
$\Sat{M}{R(x,y)}[\Subst{\Subst{s}{1}{x}}{n}{y}]$. در واقع، اگر
$n = 1$ بگیریم، داریم $\Subst{\Subst{s}{1}{x}}{n}{y} =
\Subst{s}{1}{y} = s$. چون $s(x) = 1$، $s(y) = 1$ و
$\tuple{1,1} \in \Assign{R}{M}$، پاسخ مثبت است.

برای تعیین اینکه آیا $\Sat{M}{\lexists[y][R(x,y)]}[\Subst{s}{2}{x}]$،
باید گمارش‌های !!{variable}~$\Subst{\Subst{s}{2}{x}}{n}{y}$ را بررسی
کنیم. در اینجا، برای $n = 1$، این گمارش همان
~$s_2 = \Subst{s}{2}{x}$ است، که $R(x,y)$ را ارضا نمی‌کند ($s_2(x) =
2$، $s_2(y) = 1$ و $\tuple{2,1}\notin \Assign{R}{M}$). بااین‌حال،
$\Subst{\Subst{s}{2}{x}}{3}{y} = \Subst{s_2}{3}{y}$ را در نظر بگیرید.
$\Sat{M}{R(x,y)}[\Subst{s_2}{3}{y}]$، زیرا $\tuple{2,3} \in
\Assign{R}{M}$، و بنابراین $\Sat{M}{\lexists[y][R(x,y)]}[s_2]$.

پس برای همهٔ $m \in \Domain M$، یا
$\Sat/{M}{R(a,x)}[\Subst{s}{m}{x}]$ (اگر $m = 3$، $4$) یا
$\Sat{M}{\lexists[y][R(x,y)]}[\Subst{s}{m}{x}]$ (اگر $m = 1$، $2$)؛
ازاین‌رو
\[
\Sat{M}{\lforall[x][(R(a,x) \lif \lexists[y][R(x,y)])]}[s].
\]
از سوی دیگر،
\[
\Sat/{M}{\lexists[x][(R(a,x) \land \lforall[y][R(x,y)])]}[s].
\]
فقط برای $m = 1$ و $m = 2$ داریم
$\Sat{M}{R(a,x)}[\Subst{s}{m}{x}]$. اما برای هر دو مقدار~$m$، به نوبهٔ
خود $n \in \Domain M$ای، یعنی $n = 4$، وجود دارد که
$\Sat/{M}{R(x,y)}[\Subst{\Subst{s}{m}{x}}{n}{y}]$؛ پس برای
$m = 1$ و $m = 2$،
$\Sat/{M}{\lforall[y][R(x,y)]}[\Subst{s}{m}{x}]$. در مجموع، هیچ
$m \in \Domain M$ای وجود ندارد که $\Sat{M}{R(a,x)
\land \lforall[y][R(x,y)]}[\Subst{s}{m}{x}]$.
\end{ex}

\iftag{defEx}{%
\begin{prop}\ollabel{prop:sat-ex}
$\Sat{M}{\lexists[x][!B(x)]}[s]$ اگر و تنها اگر واریانتِ $x$ای چون
$s'$ از $s$ وجود داشته باشد که $\Sat{M}{!B(x)}[s']$.
\end{prop}

\begin{proof}
  تمرین.
\end{proof}
}{}

\tagprob{defEx}
\begin{prob}
  \olref[fol][syn][sat]{prop:sat-ex} را اثبات کنید.
\end{prob}
\tagendprob

\iftag{defAll}{%
\begin{prop}\ollabel{prop:sat-all}
$\Sat{M}{\lforall[x][!B(x)]}[s]$ اگر و تنها اگر برای هر واریانتِ
$x$ چون~$s'$ از $s$، $\Sat{M}{!B(x)}[s']$.
\end{prop}

\begin{proof}
  تمرین.
\end{proof}
}{}

\tagprob{defAll}
\begin{prob}
  \olref[fol][syn][sat]{prop:sat-all} را اثبات کنید.
\end{prob}
\tagendprob

\begin{prob}
بگذارید $\Lang L = \{c, f, A\}$ دارای یک !!{constant}، یک !!{function}
یک‌جایگاهی و یک !!{predicate} دوجایگاهی باشد، و !!{structure}~$\Struct{M}$
چنین داده شود:
\begin{enumerate}
\item $\Domain M = \{1, 2, 3\}$
\item $\Assign{c}{M} = 3$
\item $\Assign{f}{M}(1) = 2, \Assign{f}{M}(2) = 3, \Assign{f}{M}(3) = 2$
\item $\Assign{A}{M} = \{\tuple{1, 2}, \tuple{2, 3}, \tuple{3, 3}\}$
\end{enumerate}
(الف) بگذارید برای همهٔ !!{variable}s~$v$، $s(v) = 1$. تعیین کنید آیا
\[
\Sat{M}{\lexists[x][(A(f(z), c) \lif \lforall[y][(A(y, x) \lor A(f(y),
      x))])]}[s]
\]
برقرار است یا نه. پاسخ خود را توضیح دهید.

(ب) ساختار و گمارش !!{variable} متفاوتی بدهید که در آن !!{formula} ارضا
نشود.
\end{prob}

\end{document}
